% Part: computability
% Chapter: computability-theory
% Section: complement-ce

\documentclass[../../../include/open-logic-section]{subfiles}

\begin{document}

\olfileid{cmp}{thy}{cmp}
\olsection[سیٹاں دا متمم]{قابلِ حساب گنتی وار فہرست بنن جوگ
  سیٹ متمم تحت بند نہیں}

من لو $A$ قابلِ حساب گنتی وار فہرست بنن جوگ اے۔
کی~$A$ دا متمم،
$\Complement{A} = \Nat \setminus A$،
ہمیشہ ایداں دا ای گنتی وار فہرست بنن جوگ
ہوندا اے؟ اگلا قضیہ تے نتیجہ دسّن گے
کہ جواب نہیں اے۔

\begin{thm}
\ollabel{thm:ce-comp}
من لو $A$ قدرتی اعداد دا کوئی وی سیٹ اے۔
تاں $A$ قابلِ حساب اے جے تے صرف جے
$A$ تے $\Complement{A}$ دوویں قابلِ حساب
گنتی وار فہرست بنن جوگ ہون۔
\end{thm}

\begin{proof}
اک پاسے دا ثبوت سوکھا اے: جے $A$ قابلِ حساب
اے تاں $\Complement{A}$ وی قابلِ حساب اے
($\Char{A} = 1 \tsub
\Char{\Complement{A}}$)؛ ایس لئی دوویں
قابلِ حساب گنتی وار فہرست بنن جوگ نیں۔

دوجے پاسے، من لو $A$ تے~$\Complement{A}$
دوویں قابلِ حساب گنتی وار فہرست بنن جوگ نیں۔
من لو $A$، $\cfind{d}$ دا دائرۂ تعریف اے،
تے $\Complement{A}$، $\cfind{e}$ دا دائرۂ
تعریف اے۔ $h$ دی تعریف ایویں کرو:
\[
h(x) = \umin{s}{(T(d,x,s) \lor T(e,x,s))}.
\]
ہور لفظاں وچ، انپٹ~$x$ اُتے $h$،
$\cfind{d}$ یا $\cfind{e}$ دی رُک جان والی
حسابی کارروائی لبھدا اے۔ ہن جے $x \in A$
اے تاں پہلی صورت وچ، تے جے $x \in
\Complement{A}$ اے تاں دوجی صورت وچ،
ایہ تلاش ضرور کامیاب ہووے گی۔ ایس لئی
$h$ اک کُل قابلِ حساب فنکشن اے۔
% PNB-CMP-018: A is the d-domain, not the e-domain in the frozen source.
ہن ہر~$x$ لئی $x \in A$ جے تے صرف
جے $T(d, x, h(x))$، یعنی جے
$\cfind{d}$ ایتھے تعریف شدہ اے۔
چونکہ $T(d, x, h(x))$ اک قابلِ حساب
تعلق اے، ایس لئی $A$ قابلِ حساب اے۔
\end{proof}

\begin{explain}
ایس گل نوں غیر رسمی حسابی لفظاں وچ سمجھنا
سوکھا اے: $A$ بارے فیصلہ کرن لئی،
انپٹ $x$ اُتے $\cfind{d}$ تے $\cfind{e}$
دیاں رُک جان والیاں حسابی کارروائیاں
لبھو۔ دوہاں وچوں اک ضرور رُکے گا؛
% PNB-CMP-019: search the d/e domains; d, not e, witnesses membership in A.
جے $\cfind{d}$ رُکے تاں $x$،~$A$
وچ اے، تے نہیں تاں $x$،~$\Complement{A}$
وچ اے۔
\end{explain}

\begin{cor}
\ollabel{cor:comp-k}
$\Complement{K_0}$ قابلِ حساب گنتی وار
فہرست بنن جوگ نہیں اے۔
\end{cor}

\begin{proof}
سانوں پتا اے کہ $K_0$ قابلِ حساب گنتی وار
فہرست بنن جوگ اے، پر قابلِ حساب نہیں۔
جے $\Complement{K_0}$ وی قابلِ حساب
گنتی وار فہرست بنن جوگ ہوندا تاں
\olref{thm:ce-comp} مطابق $K_0$ قابلِ
حساب ہوندا۔ ایہ \olref[ncp]{thm:K-0}
نال متصادم اے۔
\end{proof}

\end{document}
